% !TEX program = xelatex
% =====================================================================
%  BÁO CÁO ĐỒ ÁN TRÍ TUỆ NHÂN TẠO, HỌC PHẦN AIP202
%  Khoa Công nghệ Thông tin, Trường Đại học Kiến trúc Đà Nẵng
%
%  Build:  latexmk -xelatex do-an.tex
%  Overleaf: Menu > Compiler > XeLaTeX, rồi bấm Recompile.
% =====================================================================
% \documentclass phải nằm ở file này, không được chuyển vào settings.tex,
% vì Overleaf dò file chính bằng cách tìm file có chứa \documentclass.
\documentclass[12pt, a4paper, oneside]{book}
% !TEX root = ../do-an.tex
% =====================================================================
%  styles/settings.tex
%  Các thiết lập chung, nạp từ file chính do-an.tex.
%  SINH VIÊN KHÔNG SỬA FILE NÀY.
%
%  Lưu ý: lệnh khai báo lớp tài liệu nằm ở file do-an.tex chứ không nằm
%  ở đây. Lý do: Overleaf dò file chính bằng cách tìm chuỗi khai báo lớp
%  tài liệu trong các file. Nếu chuỗi đó xuất hiện ở file này, Overleaf
%  sẽ biên dịch nhầm file này và báo lỗi "no legal end found", vì file
%  này chỉ chứa thiết lập chứ không chứa phần thân tài liệu.
% =====================================================================

% Toàn bộ nhận diện DAU: phông, màu, bìa, hộp trình bày, header
\usepackage{styles/dau}
\dauLayout

% Giãn dòng và khoảng cách đoạn theo quy cách đồ án
\linespread{1.25}
\usepackage[skip=8pt plus 1pt, indent=15pt]{parskip}

% Ghi chú: mẫu gốc dùng gói atbegshi để vứt trang đầu, vì trang bìa của
% mẫu gốc dựng bằng minipage nên sinh ra một trang trắng. Trang bìa ở
% đây dựng bằng môi trường titlepage nên không có trang trắng, và nếu
% giữ thủ thuật đó thì chính trang bìa sẽ bị vứt mất.


% =====================================================================
%  PHẦN EM ĐIỀN: thông tin đồ án
% =====================================================================
% Loại báo cáo
\crname{BÁO CÁO ĐỒ ÁN TRÍ TUỆ NHÂN TẠO}
% Học phần
\cmonhoc{Học phần AIP202 \textperiodcentered{} 2 tín chỉ}
% Tên đề tài. Dùng \\ để ngắt dòng cho cân đối trên bìa.
\ctname{PHÂN TÍCH CẢM XÚC BÌNH LUẬN TIẾNG VIỆT\\BẰNG HỌC MÁY}
% Họ tên sinh viên và mã số sinh viên
\cstuname{Nguyễn Văn A}
\cmssv{2022xxxxx}
% Lớp
\clopsv{CNTT 22A}
% Ngành
\csCouncil{Công nghệ thông tin}
% Giảng viên hướng dẫn
\csSupervise{TS. Đỗ Phúc Hảo}
% Năm học
\cttime{2025--2026}
% Khoa quản lý
\csdeptname{KHOA CÔNG NGHỆ THÔNG TIN}
% Nơi thực hiện
\csplace{Đà Nẵng}

\begin{document}

% Bìa trang trọng có khung kép. Muốn dùng bìa hiện đại theo bộ giáo
% trình của khoa thì thay dòng dưới bằng \coverpagehiendai
\coverpage

% ---------------------------------------------------------------------
%  PHẦN ĐẦU: đánh số trang bằng số La Mã thường
% ---------------------------------------------------------------------
\frontmatter
\dauPageStyle

% ---------------------------------------------------------------------
%  LỜI CAM ĐOAN
% ---------------------------------------------------------------------
\begin{declaration}

\begin{chuy}
\goiy{giữ nguyên tinh thần đoạn dưới, chỉ thay tên giảng viên hướng dẫn và
tên em. Đây là cam kết có giá trị pháp lý về liêm chính học thuật, đừng viết
qua loa. Xoá hộp gợi ý này sau khi sửa xong.}
\end{chuy}

Tôi xin cam đoan đồ án này là công trình do chính tôi thực hiện dưới sự
hướng dẫn của \dauGVHD. Các nội dung nghiên cứu, số liệu và kết quả trình
bày trong đồ án là trung thực và chưa từng được công bố trong bất kỳ công
trình nào khác.

Mọi số liệu, hình ảnh, mã nguồn và kết quả của tác giả khác được sử dụng
trong đồ án đều được trích dẫn nguồn gốc rõ ràng trong phần tài liệu tham
khảo. Toàn bộ phần mềm và thư viện được dùng đều có giấy phép cho phép sử
dụng. Những phần có sử dụng công cụ trí tuệ nhân tạo hỗ trợ đã được nêu rõ
trong báo cáo.

Nếu phát hiện có bất kỳ sự gian lận nào, tôi xin chịu hoàn toàn trách nhiệm
trước Khoa và Nhà trường.

\vspace{1.4em}
\hfill\okyten{\dauSinhVien}{Sinh viên thực hiện}

\end{declaration}

% ---------------------------------------------------------------------
%  LỜI CẢM ƠN
% ---------------------------------------------------------------------
\begin{acknowledgments}

\begin{chuy}
\goiy{viết một đoạn khoảng năm đến tám dòng, cảm ơn giảng viên hướng dẫn,
khoa, gia đình và bạn bè. Nếu đồ án có dùng dữ liệu, mã nguồn hay hạ tầng
tính toán của đơn vị nào thì ghi lời cảm ơn ở đây luôn.}
\end{chuy}

\begin{mau}
Em xin gửi lời cảm ơn chân thành đến \dauGVHD, người đã tận tình định hướng
đề tài, góp ý và theo sát tiến độ của em trong suốt quá trình thực hiện đồ
án. Em cũng xin cảm ơn quý thầy cô Khoa Công nghệ Thông tin, Trường Đại học
Kiến trúc Đà Nẵng đã trang bị cho em nền tảng kiến thức để hoàn thành công
việc này. Cuối cùng, em xin cảm ơn gia đình và bạn bè đã luôn động viên em
trong thời gian học tập.
\end{mau}

\vspace{1.4em}
\hfill\okyten{\dauSinhVien}{Sinh viên thực hiện}

\end{acknowledgments}

% ---------------------------------------------------------------------
%  TÓM TẮT
% ---------------------------------------------------------------------
\begin{abstract}

\begin{chuy}
\goiy{tóm tắt khoảng 200 đến 300 từ, đi theo năm ý và mỗi ý một đến hai câu:
bối cảnh, vấn đề cần giải, cách em giải, kết quả chính kèm con số, và ý
nghĩa. Đây là phần hội đồng đọc đầu tiên, nên viết sau cùng khi đã có kết
quả thật. Không đưa trích dẫn và không đưa nội dung chưa có trong báo cáo.}
\end{chuy}

\begin{mau}
Nhận dạng loài vật qua ảnh là bài toán có nhiều ứng dụng trong giám sát đa
dạng sinh học và giáo dục. Thách thức nằm ở chỗ ảnh chụp ngoài thực địa có
nền phức tạp, điều kiện ánh sáng thay đổi và nhiều loài trông gần giống
nhau. Đồ án xây dựng một hệ thống nhận dạng loài vật dựa trên mạng nơ-ron
tích chập, kết hợp học chuyển giao từ mô hình đã tiền huấn luyện và các phép
tăng cường dữ liệu phù hợp với đặc thù ảnh thực địa. Hệ thống được huấn
luyện và đánh giá trên tập dữ liệu gồm 12.000 ảnh thuộc 30 loài, đạt độ
chính xác 91,2\% trên tập kiểm thử độc lập, cao hơn 6,4 điểm phần trăm so
với mô hình huấn luyện từ đầu. Toàn bộ mô hình được triển khai thành một ứng
dụng web cho phép người dùng tải ảnh lên và nhận kết quả kèm độ tin cậy
trong dưới một giây. Kết quả cho thấy học chuyển giao là lựa chọn hợp lý khi
dữ liệu huấn luyện ở quy mô vừa, và hệ thống đủ nhẹ để chạy trên máy tính
cá nhân thông thường.

\vspace{0.8em}
\noindent\textbf{Từ khoá:} mạng nơ-ron tích chập, học chuyển giao, nhận dạng
ảnh, thị giác máy tính.
\end{mau}

\end{abstract}


\tableofcontents
\listoffigures
\listoftables
% Bỏ dòng dưới nếu đồ án của em không có giải thuật nào
\listofalgorithms

% ---------------------------------------------------------------------
%  THÂN BÁO CÁO: sáu chương theo quy cách của học phần
% ---------------------------------------------------------------------
\mainmatter
\dauPageStyle

% ---------------------------------------------------------------------
%  CHƯƠNG 1 — MỞ ĐẦU
% ---------------------------------------------------------------------
\chapter{Mở đầu}
\label{ch:modau}

\begin{muctieu}
Chương mở đầu trình bày lý do chọn đề tài, mục tiêu cần đạt, phạm vi thực
hiện và bố cục báo cáo, giúp người đọc nắm được bức tranh tổng quan về đồ án.
\end{muctieu}

\section{Lý do chọn đề tài}

\begin{chuy}
\goiy{nêu bối cảnh thực tế của bài toán, vì sao nó đáng giải quyết, và vì
sao em chọn cách tiếp cận bằng trí tuệ nhân tạo. Có thể dẫn một đến hai số
liệu hoặc một tình huống minh hoạ.}
\end{chuy}

\begin{mau}
Trong bối cảnh thương mại điện tử và mạng xã hội phát triển mạnh, mỗi ngày
có hàng triệu bình luận của người dùng được tạo ra. Việc đọc và tổng hợp
thủ công các bình luận này tốn nhiều công sức và không kịp thời. Bài toán
phân tích cảm xúc \en{(sentiment analysis)} cho phép tự động xác định một
bình luận mang sắc thái tích cực, tiêu cực hay trung tính, từ đó hỗ trợ
doanh nghiệp lắng nghe khách hàng nhanh hơn. Đề tài này được chọn vì vừa có
tính thực tiễn cao, vừa giúp em vận dụng các kiến thức về xử lý ngôn ngữ tự
nhiên và học máy đã tích luỹ.
\end{mau}

\section{Mục tiêu của đồ án}

\begin{chuy}
\goiy{liệt kê ba đến năm mục tiêu cụ thể, đo lường được. Tránh viết chung
chung kiểu "tìm hiểu về AI". Hãy nêu rõ sản phẩm và mức độ hoàn thành mong
muốn, vì Chương~\ref{ch:ketqua} sẽ đối chiếu ngược lại từng mục tiêu này.}
\end{chuy}

\begin{mau}
Đồ án hướng tới các mục tiêu sau:
\begin{itemize}[leftmargin=1.6em, itemsep=2pt]
  \item Xây dựng được một mô hình phân loại cảm xúc cho bình luận tiếng Việt
        với độ chính xác từ 80\% trở lên trên tập kiểm thử.
  \item Cài đặt quy trình tiền xử lý dữ liệu văn bản tiếng Việt hoàn chỉnh.
  \item Phát triển giao diện web cho phép người dùng nhập câu và nhận kết
        quả dự đoán.
  \item So sánh và đánh giá ít nhất hai mô hình học máy khác nhau.
\end{itemize}
\end{mau}

\section{Phạm vi và đối tượng}

\begin{chuy}
\goiy{giới hạn rõ những gì đồ án làm và không làm: loại dữ liệu, ngôn ngữ,
số lượng nhãn, nền tảng triển khai. Giới hạn phạm vi sớm là cách chắc chắn
nhất để kịp hoàn thành trong mười hai tuần.}
\end{chuy}

\begin{mau}
Đồ án tập trung vào phân loại cảm xúc ở mức câu cho văn bản tiếng Việt, với
ba nhãn tích cực, tiêu cực và trung tính. Dữ liệu là bình luận sản phẩm thu
thập công khai. Đồ án không xử lý dữ liệu đa ngôn ngữ, không phát hiện châm
biếm và chưa triển khai ở quy mô sản xuất.
\end{mau}

\begin{luuy}
Lỗi phổ biến nhất của đồ án là chọn phạm vi quá rộng rồi bỏ dở. Một bài toán
hẹp làm tới nơi tới chốn luôn được chấm cao hơn một ý tưởng lớn còn dang dở.
\end{luuy}

\section{Bố cục báo cáo}

\begin{chuy}
\goiy{giữ nguyên đoạn dưới, chỉ sửa nếu em thêm hoặc bớt chương.}
\end{chuy}

Phần còn lại của báo cáo được tổ chức như sau. Chương~\ref{ch:coso} trình
bày cơ sở lý thuyết và công nghệ sử dụng. Chương~\ref{ch:thietke} phân tích
bài toán và thiết kế giải pháp. Chương~\ref{ch:caidat} mô tả quá trình cài
đặt và thử nghiệm. Chương~\ref{ch:ketqua} trình bày kết quả và đánh giá.
Chương~\ref{ch:ketluan} tổng kết và nêu hướng phát triển.

% ---------------------------------------------------------------------
%  CHƯƠNG 2 — CƠ SỞ LÝ THUYẾT VÀ CÔNG NGHỆ
% ---------------------------------------------------------------------
\chapter{Cơ sở lý thuyết và công nghệ}
\label{ch:coso}

\begin{muctieu}
Chương này trình bày các kiến thức nền tảng gồm thuật toán và mô hình trí
tuệ nhân tạo được dùng, cùng danh mục công cụ và thư viện để xây dựng hệ thống.
\end{muctieu}

\section{Kiến thức nền tảng}

\begin{chuy}
\goiy{trình bày ngắn gọn lý thuyết của thuật toán hay mô hình mà em dùng,
ví dụ phân lớp, mạng nơ-ron, cây quyết định, thuật toán tìm kiếm heuristic.
Chỉ nêu phần liên quan trực tiếp tới đồ án, kèm công thức nếu cần. Đây không
phải chương chép lại giáo trình.}
\end{chuy}

\begin{dinhnghia}[Phân lớp văn bản]
Phân lớp văn bản \en{(text classification)} là bài toán gán cho một đoạn văn
bản $x$ một nhãn $y$ trong tập nhãn hữu hạn $\mathcal{Y}$. Bộ phân lớp học
một hàm $f: x \mapsto y$ từ dữ liệu huấn luyện đã gán nhãn.
\end{dinhnghia}

\begin{mau}
Một hướng tiếp cận phổ biến là biểu diễn văn bản theo mô hình túi từ
\en{(bag of words)} kết hợp trọng số TF-IDF, sau đó dùng bộ phân lớp tuyến
tính. Trọng số TF-IDF của từ $t$ trong văn bản $d$ được tính theo công thức
\begin{equation}
\mathrm{tfidf}(t, d) = \mathrm{tf}(t, d) \cdot \log\frac{N}{1 + \mathrm{df}(t)},
\label{eq:tfidf}
\end{equation}
trong đó $\mathrm{tf}(t,d)$ là tần suất từ $t$ trong văn bản $d$, $N$ là
tổng số văn bản và $\mathrm{df}(t)$ là số văn bản có chứa từ $t$. Thừa số
logarit trong~\eqref{eq:tfidf} làm giảm trọng số của những từ xuất hiện ở
hầu hết văn bản, vốn ít giá trị phân biệt.
\end{mau}

\begin{ghinho}
Với dữ liệu vừa và nhỏ, các mô hình tuyến tính như hồi quy logistic hay máy
vec-tơ hỗ trợ thường cho kết quả tốt mà chi phí huấn luyện thấp, rất hợp với
phạm vi một đồ án học phần.
\end{ghinho}

\section{Công cụ và thư viện}

\begin{chuy}
\goiy{liệt kê ngôn ngữ, framework và thư viện chính, ghi rõ phiên bản và vai
trò từng thành phần. Bảng phiên bản chính là thứ giúp người khác chạy lại
được sản phẩm của em.}
\end{chuy}

\begin{table}[H]
\centering
\caption{Công cụ và thư viện sử dụng trong đồ án.}
\label{tab:congcu}
\begin{tabular*}{\textwidth}{@{\extracolsep{\fill}}>{\bfseries\color{daunavy}}l >{\color{daumaroon}}l l}
\toprule
\rowcolor{gray100}\bfseries\color{daunavy}Thành phần & \bfseries\color{daumaroon}Phiên bản & \bfseries\color{daunavy}Vai trò \\
\midrule
Python       & 3.11 & Ngôn ngữ lập trình chính \\
scikit-learn & 1.4  & Huấn luyện và đánh giá mô hình học máy \\
pandas       & 2.2  & Đọc và xử lý dữ liệu dạng bảng \\
underthesea  & 6.x  & Tách từ tiếng Việt \\
Flask        & 3.0  & Xây dựng giao diện web demo \\
\bottomrule
\end{tabular*}
\end{table}

\section{Các công trình liên quan}

\begin{chuy}
\goiy{tóm tắt hai đến ba cách làm sẵn có cho bài toán của em, kèm trích dẫn
bằng lệnh \code{cite} với khoá khai báo trong file \code{refs.bib}. Nêu rõ em
kế thừa gì và làm khác gì. Mục này không bắt buộc nhưng làm báo cáo có chiều
sâu hơn hẳn.}
\end{chuy}

\begin{mau}
Các nghiên cứu về phân tích cảm xúc tiếng Việt thường đi theo hai hướng.
Hướng thứ nhất dùng đặc trưng thủ công kết hợp mô hình tuyến tính, ưu điểm
là nhẹ và dễ giải thích~\cite{HoangKiem}. Hướng thứ hai dùng mô hình học sâu
tiền huấn luyện cho độ chính xác cao hơn nhưng đòi hỏi tài nguyên tính toán
lớn~\cite{Krizhevsky2012}. Đồ án này
chọn hướng thứ nhất làm cơ sở, đồng thời so sánh với một mô hình mạnh hơn để
thấy rõ đánh đổi giữa độ chính xác và chi phí.
\end{mau}

% ---------------------------------------------------------------------
%  CHƯƠNG 3 — PHÂN TÍCH VÀ THIẾT KẾ
% ---------------------------------------------------------------------
\chapter{Phân tích và thiết kế}
\label{ch:thietke}

\begin{muctieu}
Chương này phân tích yêu cầu bài toán, đề xuất giải pháp, thiết kế kiến trúc
hệ thống, thiết kế dữ liệu và thiết kế thành phần tri thức hoặc mô hình.
\end{muctieu}

\section{Phân tích bài toán và yêu cầu}

\begin{chuy}
\goiy{mô tả đầu vào, đầu ra của hệ thống và các yêu cầu chức năng, phi chức
năng. Có thể dùng kịch bản sử dụng \en{(use case)} để làm rõ. Hãy nêu luôn
tiêu chí thế nào là hệ thống chạy đúng.}
\end{chuy}

\begin{mau}
Hệ thống nhận đầu vào là một câu bình luận tiếng Việt và trả về nhãn cảm xúc
dự đoán kèm độ tin cậy. Yêu cầu chức năng gồm tiền xử lý văn bản, dự đoán
cảm xúc và hiển thị kết quả. Yêu cầu phi chức năng gồm thời gian phản hồi
dưới một giây cho mỗi câu và giao diện dùng được trên trình duyệt phổ thông.
\end{mau}

\section{Thiết kế kiến trúc hệ thống}

\begin{chuy}
\goiy{vẽ sơ đồ kiến trúc tổng thể gồm các khối và luồng dữ liệu. Em có thể
sửa lại sơ đồ TikZ mẫu bên dưới, hoặc chèn ảnh sơ đồ qua lệnh
\code{includegraphics}. Mọi hình đều phải được nhắc tới trong thân bài.}
\end{chuy}

Kiến trúc hệ thống được trình bày ở Hình~\ref{fig:kientruc}. Dữ liệu đi tuần
tự qua bốn khối, mỗi khối là một mô-đun độc lập trong mã nguồn.

\begin{figure}[H]
\centering
\resizebox{\textwidth}{!}{%
\begin{tikzpicture}[node distance=12mm and 14mm]
  \node[objbox] (in) {Bình luận\\ đầu vào};
  \node[clsbox, right=of in]  (pre) {Tiền xử lý\\ (tách từ, làm sạch)};
  \node[clsbox, right=of pre] (vec) {Vec-tơ hoá\\ TF-IDF};
  \node[clsbox, right=of vec] (mod) {Mô hình\\ phân lớp};
  \node[objbox, right=of mod] (out) {Nhãn cảm xúc\\ và độ tin cậy};
  \draw[flowarr] (in)  -- (pre);
  \draw[flowarr] (pre) -- (vec);
  \draw[flowarr] (vec) -- (mod);
  \draw[flowarr] (mod) -- (out);
\end{tikzpicture}}
\caption{Sơ đồ kiến trúc tổng thể của hệ thống.}
\label{fig:kientruc}
\end{figure}

\section{Thiết kế dữ liệu}

\begin{chuy}
\goiy{mô tả nguồn dữ liệu, số lượng mẫu, cấu trúc các trường và cách chia
tập huấn luyện với tập kiểm thử. Nếu có cơ sở dữ liệu thì mô tả lược đồ
bảng. Ghi rõ xuất xứ dữ liệu, đây là câu hỏi giảng viên hay hỏi khi bảo vệ.}
\end{chuy}

\begin{mau}
Bộ dữ liệu gồm 5.000 bình luận đã gán nhãn, chia theo tỉ lệ 80\% huấn luyện
và 20\% kiểm thử, chia phân tầng theo nhãn để giữ nguyên tỉ lệ các lớp. Cấu
trúc mỗi bản ghi được mô tả ở Bảng~\ref{tab:dulieu}.
\end{mau}

\begin{table}[H]
\centering
\caption{Cấu trúc bản ghi dữ liệu.}
\label{tab:dulieu}
\begin{tabular*}{\textwidth}{@{\extracolsep{\fill}}>{\bfseries\color{daunavy}}l >{\color{daumaroon}}l l}
\toprule
\rowcolor{gray100}\bfseries\color{daunavy}Trường & \bfseries\color{daumaroon}Kiểu & \bfseries\color{daunavy}Mô tả \\
\midrule
id      & integer & Khoá định danh bản ghi \\
content & text    & Nội dung bình luận \\
label   & text    & Nhãn cảm xúc: tích cực, tiêu cực, trung tính \\
\bottomrule
\end{tabular*}
\end{table}

\section{Thiết kế mô hình hoặc cơ sở tri thức}

\begin{chuy}
\goiy{đây là phần lõi của đồ án và tương ứng chuẩn đầu ra CLO3. Nếu em làm
hệ chuyên gia thì mô tả tập luật và cơ chế suy diễn. Nếu em làm học máy thì
mô tả kiến trúc mô hình, siêu tham số và cách chọn chúng. Đừng chỉ viết em
dùng thư viện nào, hãy nói rõ vì sao chọn cấu hình đó.}
\end{chuy}

\begin{mau}
Mô hình được chọn là hồi quy logistic trên đặc trưng TF-IDF với n-gram bậc
một và bậc hai, ngưỡng tần suất tối thiểu \code{min\_df} bằng 2 để loại bỏ
những từ chỉ xuất hiện một lần. Siêu tham số được chọn bằng tìm kiếm lưới
kết hợp kiểm định chéo năm phần trên tập huấn luyện, không đụng tới tập kiểm
thử trong quá trình chọn tham số.
\end{mau}

\begin{luuy}
Không được dùng tập kiểm thử để chọn siêu tham số. Làm vậy thì con số độ
chính xác báo cáo sẽ cao hơn thực tế, và em sẽ không giải thích được khi
hội đồng hỏi tới.
\end{luuy}


%% ---------------------------------------------------------------------
%  CHƯƠNG 4 — CÀI ĐẶT VÀ THỬ NGHIỆM
% ---------------------------------------------------------------------
\chapter{Cài đặt và thử nghiệm}
\label{ch:caidat}

\begin{muctieu}
Chương này mô tả môi trường triển khai, các thành phần chính kèm mã nguồn
minh hoạ, và giao diện demo của hệ thống.
\end{muctieu}

\section{Môi trường triển khai}

\begin{chuy}
\goiy{ghi rõ cấu hình phần cứng, hệ điều hành và phiên bản phần mềm để người
khác tái lập được kết quả. Kèm cả đường dẫn kho \code{git} chứa mã nguồn.}
\end{chuy}

\begin{mau}
Hệ thống được phát triển và thử nghiệm trên máy tính cá nhân với CPU Intel
Core i5, RAM 8~GB, hệ điều hành Windows 11 và Python 3.11 cùng các thư viện
đã nêu ở Bảng~\ref{tab:congcu}. Mã nguồn được quản lý bằng \code{git} và đặt
tại kho \url{https://github.com/tencuaem/ten-do-an}.
\end{mau}

\section{Các thành phần chính}

\begin{chuy}
\goiy{trình bày những đoạn mã quan trọng nhất như tiền xử lý, huấn luyện,
dự đoán. Chỉ đưa phần cốt lõi, không dán toàn bộ mã nguồn vào thân báo cáo;
mã dài để ở phụ lục. Comment trong code viết bằng tiếng Anh theo quy ước của
khoa.}
\end{chuy}

Mã nguồn~\ref{lst:train} minh hoạ bước vec-tơ hoá và huấn luyện mô hình.

\begin{lstlisting}[language=Python, caption={Huấn luyện mô hình phân lớp cảm xúc.}, label={lst:train}]
from sklearn.feature_extraction.text import TfidfVectorizer
from sklearn.linear_model import LogisticRegression
from sklearn.pipeline import Pipeline

# Build a pipeline: TF-IDF features feed a logistic regression classifier
model = Pipeline([
    ("tfidf", TfidfVectorizer(ngram_range=(1, 2), min_df=2)),
    ("clf", LogisticRegression(max_iter=1000)),
])

# Fit the model on the training split only
model.fit(X_train, y_train)
\end{lstlisting}

Mã nguồn~\ref{lst:predict} minh hoạ hàm dự đoán dùng cho giao diện web.

\begin{lstlisting}[language=Python, caption={Hàm dự đoán cảm xúc cho một câu.}, label={lst:predict}]
def predict_sentiment(text):
    # Return the predicted label together with its confidence score
    label = model.predict([text])[0]
    confidence = model.predict_proba([text]).max()
    return label, round(float(confidence), 3)
\end{lstlisting}

\section{Giao diện và demo}

\begin{chuy}
\goiy{chèn ảnh chụp màn hình giao diện chính. Đặt file ảnh vào thư mục
\code{hinh/}, rồi bỏ dấu phần trăm ở dòng \code{includegraphics} bên dưới và
xoá khung giữ chỗ.}
\end{chuy}

\begin{figure}[H]
\centering
% Bỏ dấu % ở dòng dưới và thay ten-anh.png bằng ảnh thật của em
% \includegraphics[width=0.8\textwidth]{ten-anh.png}
\fbox{\begin{minipage}[c][5cm][c]{0.8\textwidth}\centering
  \itshape\color{gray500}[Chèn ảnh chụp màn hình giao diện demo tại đây]
\end{minipage}}
\caption{Giao diện web cho phép nhập câu và xem kết quả dự đoán.}
\label{fig:demo}
\end{figure}

\section{Kiểm thử}

\begin{chuy}
\goiy{nêu cách em kiểm tra hệ thống chạy đúng: các ca kiểm thử đã chạy, dữ
liệu biên đã thử, lỗi đã phát hiện và cách khắc phục. Một bảng ca kiểm thử
ngắn cũng đủ và rất được đánh giá cao.}
\end{chuy}

\begin{mau}
Hệ thống được kiểm thử với ba nhóm ca. Nhóm một là câu ngắn rõ sắc thái,
nhóm hai là câu dài nhiều mệnh đề, nhóm ba là dữ liệu biên gồm chuỗi rỗng,
chuỗi chỉ có ký tự đặc biệt và chuỗi rất dài. Ở nhóm ba, hệ thống ban đầu
báo lỗi với chuỗi rỗng; lỗi được khắc phục bằng cách kiểm tra đầu vào trước
khi gọi mô hình.
\end{mau}

%% ---------------------------------------------------------------------
%  CHƯƠNG 5 — KẾT QUẢ VÀ ĐÁNH GIÁ
% ---------------------------------------------------------------------
\chapter{Kết quả và đánh giá}
\label{ch:ketqua}

\begin{muctieu}
Chương này trình bày kết quả đạt được, các độ đo đánh giá định lượng, phần
bàn luận và những hạn chế còn tồn tại.
\end{muctieu}

\section{Kết quả đạt được}

\begin{chuy}
\goiy{đối chiếu trực tiếp từng mục tiêu đã nêu ở Chương~\ref{ch:modau} để
người đọc thấy mức độ hoàn thành. Mục tiêu nào chưa đạt thì nói thẳng, đừng
lờ đi.}
\end{chuy}

\begin{mau}
Hệ thống đã hoàn thành các mục tiêu đề ra. Mô hình phân loại đạt độ chính
xác vượt ngưỡng 80\% trên tập kiểm thử, quy trình tiền xử lý hoạt động ổn
định và giao diện web cho phép dự đoán theo thời gian thực.
\end{mau}

\section{Độ đo đánh giá}

\begin{chuy}
\goiy{trình bày các độ đo phù hợp với bài toán, ví dụ độ chính xác,
precision, recall, F1 cho bài toán phân lớp, hoặc MAE và RMSE cho bài toán
dự báo. Ghi rõ số liệu đo trên tập nào và với bao nhiêu mẫu.}
\end{chuy}

Bảng~\ref{tab:ketqua} so sánh hai mô hình trên cùng tập kiểm thử gồm 1.000 mẫu.

\begin{table}[H]
\centering
\caption{So sánh hiệu năng các mô hình trên tập kiểm thử.}
\label{tab:ketqua}
\begin{tabular*}{\textwidth}{@{\extracolsep{\fill}}>{\bfseries\color{daunavy}}l cccc}
\toprule
\rowcolor{gray100}\bfseries\color{daunavy}Mô hình & \bfseries\color{daumaroon}Accuracy & \bfseries\color{daumaroon}Precision & \bfseries\color{daumaroon}Recall & \bfseries\color{daumaroon}F1 \\
\midrule
Hồi quy logistic  & 0,84 & 0,83 & 0,82 & 0,83 \\
Máy vec-tơ hỗ trợ & 0,86 & 0,85 & 0,84 & 0,85 \\
\bottomrule
\end{tabular*}
\end{table}

\section{Bàn luận}

\begin{chuy}
\goiy{giải thích vì sao kết quả lại như vậy, chứ không chỉ đọc lại bảng số.
Mô hình sai ở những ca nào và vì sao. Đây là phần phân biệt báo cáo khá với
báo cáo tốt.}
\end{chuy}

\begin{mau}
Máy vec-tơ hỗ trợ nhỉnh hơn hồi quy logistic khoảng hai điểm phần trăm ở mọi
độ đo. Khoảng cách này không lớn và phải trả giá bằng thời gian huấn luyện
lâu hơn. Phân tích các ca dự đoán sai cho thấy phần lớn lỗi tập trung ở nhãn
trung tính, nơi ranh giới ngữ nghĩa với hai nhãn còn lại là mờ nhất.
\end{mau}

\section{Hạn chế}

\begin{chuy}
\goiy{thẳng thắn nêu những điểm hệ thống chưa làm tốt: dữ liệu thiếu, lỗi
với câu phức tạp, chưa tối ưu tốc độ. Phần này thể hiện tư duy phản biện và
luôn được hội đồng ghi nhận.}
\end{chuy}

\begin{luuy}
Mô hình còn nhầm lẫn với các câu mang sắc thái châm biếm và với bình luận
chứa nhiều từ viết tắt hoặc tiếng lóng. Tập dữ liệu chưa cân bằng giữa các
nhãn cũng làm giảm độ chính xác ở nhãn trung tính.
\end{luuy}

%% ---------------------------------------------------------------------
%  CHƯƠNG 6 — KẾT LUẬN VÀ HƯỚNG PHÁT TRIỂN
% ---------------------------------------------------------------------
\chapter{Kết luận và hướng phát triển}
\label{ch:ketluan}

\begin{muctieu}
Chương này tổng kết những gì đồ án đã làm được, rút ra bài học và đề xuất
hướng phát triển tiếp theo.
\end{muctieu}

\section{Kết luận}

\begin{chuy}
\goiy{tóm tắt lại bài toán, giải pháp và kết quả chính trong vài câu. Nhấn
mạnh giá trị thực tiễn và kiến thức em thu nhận được. Không đưa thông tin
mới chưa từng xuất hiện ở các chương trước.}
\end{chuy}

\begin{mau}
Đồ án đã xây dựng thành công một hệ thống phân tích cảm xúc bình luận tiếng
Việt, đi trọn từ khâu xử lý dữ liệu, huấn luyện mô hình đến triển khai giao
diện web. Qua đó em đã vận dụng được kiến thức về xử lý ngôn ngữ tự nhiên và
học máy vào một bài toán thực tế, đồng thời làm quen với quy trình phát
triển phần mềm có quản lý phiên bản.
\end{mau}

\section{Bài học rút ra}

\begin{chuy}
\goiy{đây là phần phản tư, tương ứng chuẩn đầu ra CLO4. Nêu vai trò và giới
hạn của trí tuệ nhân tạo trong chính bài toán của em, cùng điều em sẽ làm
khác đi nếu bắt đầu lại. Viết thật, đừng viết sáo.}
\end{chuy}

\begin{mau}
Bài học lớn nhất là chất lượng dữ liệu quyết định nhiều hơn việc chọn mô
hình. Phần lớn thời gian của đồ án dành cho làm sạch và gán nhãn dữ liệu,
trong khi việc đổi mô hình chỉ đem lại vài điểm phần trăm. Nếu làm lại, em
sẽ dành hai tuần đầu để chuẩn hoá dữ liệu trước khi viết dòng mã huấn luyện
đầu tiên.
\end{mau}

\section{Hướng phát triển}

\begin{chuy}
\goiy{đề xuất ba đến bốn hướng mở rộng khả thi: dữ liệu lớn hơn, mô hình
mạnh hơn, thêm tính năng, triển khai thực tế. Mỗi hướng nói rõ nó khắc phục
hạn chế nào đã nêu ở Chương~\ref{ch:ketqua}.}
\end{chuy}

\begin{mau}
Một số hướng phát triển tiếp theo:
\begin{itemize}[leftmargin=1.6em, itemsep=2pt]
  \item Mở rộng tập dữ liệu và cân bằng số mẫu giữa các nhãn, nhằm khắc phục
        độ chính xác thấp ở nhãn trung tính.
  \item Thử nghiệm mô hình ngôn ngữ tiền huấn luyện như PhoBERT để nâng độ
        chính xác trên câu dài nhiều mệnh đề.
  \item Bổ sung phát hiện câu châm biếm và chuẩn hoá tiếng lóng.
  \item Triển khai hệ thống lên nền tảng đám mây để phục vụ nhiều người dùng
        đồng thời.
\end{itemize}
\end{mau}


% Chương tra cứu quy cách trình bày. XOÁ DÒNG DƯỚI khi nộp bản chính
% thức, vì đây chỉ là sổ tra cứu chứ không phải một chương của báo cáo.


%% ---------------------------------------------------------------------
%  CHƯƠNG 2 — QUY CÁCH TRÌNH BÀY
%  Chương này đóng vai trò sổ tra cứu: mỗi mục là một mẫu em copy rồi sửa.
%  Khi viết đồ án thật, em thay cả chương này bằng nội dung của mình.
% ---------------------------------------------------------------------
\chapter{Quy cách trình bày trong báo cáo}
\label{ch:soanthao}

\begin{muctieu}
Chương này gom sẵn mẫu cho mọi thành phần em cần dùng: hình, bảng, công
thức, mã nguồn, giải thuật, trích dẫn, chú thích và các hộp nhấn mạnh. Mỗi
mục có đoạn mã mẫu, em copy rồi thay nội dung.
\end{muctieu}

\begin{luuy}
Khi bắt đầu viết đồ án thật, hãy xoá cả chương này. Nó chỉ là sổ tra cứu,
không phải một chương của báo cáo.
\end{luuy}

\section{Chèn hình ảnh}

Đặt file ảnh vào thư mục \code{figs/} rồi dùng mẫu dưới đây. Mọi hình đều
phải có chú thích, có nhãn, và phải được nhắc tới trong thân bài như
Hình~\ref{fig:vidu}.

\begin{figure}[H]
\centering
\includegraphics[width=0.72\textwidth]{dau-campus.png}
\caption{Một ví dụ về chèn hình ảnh từ file.}
\label{fig:vidu}
\end{figure}

Nếu hình là sơ đồ do em tự vẽ thì nên vẽ thẳng bằng TikZ như
Hình~\ref{fig:luongxuly}, vì sơ đồ vẽ bằng TikZ luôn nét khi phóng to và
dùng đúng bộ màu của trường.

\begin{figure}[H]
\centering
% \resizebox ép sơ đồ vừa đúng bề ngang trang, tránh tràn lề
\resizebox{\textwidth}{!}{%
\begin{tikzpicture}[node distance=8mm and 12mm]
  \node[io] (in) {Ảnh đầu vào};
  \node[process, right=of in] (pre) {Tiền xử lý};
  \node[decision, right=of pre] (chk) {Đủ chất lượng?};
  \node[process, right=of chk] (mod) {Mô hình CNN};
  \node[startstop, below=of chk] (loai) {Loại bỏ};
  \draw[flowarr] (in)  -- (pre);
  \draw[flowarr] (pre) -- (chk);
  \draw[flowarr] (chk) -- node[above, font=\scriptsize] {có} (mod);
  \draw[flowarr] (chk) -- node[right, font=\scriptsize] {không} (loai);
\end{tikzpicture}}
\caption{Một ví dụ sơ đồ luồng xử lý vẽ bằng TikZ.}
\label{fig:luongxuly}
\end{figure}

\section{Chèn bảng}

Bảng dùng chú thích đặt phía trên, đúng quy cách trình bày khoa học. Dùng
\code{booktabs} với ba loại đường kẻ \code{toprule}, \code{midrule},
\code{bottomrule}, tránh kẻ dọc cho mọi ô.

\begin{table}[H]
\centering
\caption{Một ví dụ về bảng so sánh kết quả.}
\label{tab:sosanh}
\begin{tabular*}{\textwidth}{@{\extracolsep{\fill}}>{\bfseries\color{daunavy}}l cccc}
\toprule
\rowcolor{gray100}\bfseries\color{daunavy}Mô hình & \bfseries\color{daumaroon}Accuracy & \bfseries\color{daumaroon}Precision & \bfseries\color{daumaroon}Recall & \bfseries\color{daumaroon}F1 \\
\midrule
Huấn luyện từ đầu   & 0,848 & 0,841 & 0,836 & 0,838 \\
Học chuyển giao     & 0,912 & 0,907 & 0,903 & 0,905 \\
\bottomrule
\end{tabular*}
\end{table}

\begin{luuy}
Bảng rộng quá lề là lỗi hay gặp nhất. Nếu bảng tràn, hãy giảm cỡ chữ bằng
\code{small} hoặc \code{footnotesize}, hoặc gộp bớt cột. Đừng để LaTeX tự
tràn rồi bỏ qua.
\end{luuy}

\section{Chèn công thức toán học}

Công thức đặt giữa dòng và có đánh số thì dùng môi trường \code{equation},
như Công thức~\eqref{eq:softmax}. Công thức ngắn chèn trong câu thì viết
giữa hai dấu đô la, ví dụ $y = f(x)$.

\begin{equation}
p(y = k \mid \mathbf{x}) = \frac{\exp(z_k)}{\sum_{j=1}^{K} \exp(z_j)},
\label{eq:softmax}
\end{equation}

trong đó $z_k$ là điểm số thô mà mạng gán cho lớp $k$, và $K$ là số lớp. Mọi
ký hiệu xuất hiện trong công thức đều phải được giải thích ngay sau đó, như
câu này đang làm với \eqref{eq:softmax}.

\begin{ghinho}
Công thức là một phần của câu văn, không phải một khối trôi nổi. Trước công
thức phải có câu dẫn vào, sau công thức phải có câu giải thích ký hiệu hoặc
ý nghĩa. Công thức đứng một mình không ai đọc.
\end{ghinho}

\section{Chèn mã nguồn}

Mẫu này dùng gói \code{listings} thay cho \code{minted} của bản gốc, vì
\code{listings} không cần bật \code{-shell-escape} nên build được ngay trên
Overleaf lẫn máy cá nhân.

Cách thứ nhất là gõ trực tiếp, kết quả như Mã nguồn~\ref{lst:tructiep}.

\begin{lstlisting}[language=Python, caption={Chèn mã nguồn trực tiếp vào báo cáo.}, label={lst:tructiep}]
import torch.nn as nn

# A small convolutional block: conv, batch norm, activation
def conv_block(in_ch, out_ch):
    return nn.Sequential(
        nn.Conv2d(in_ch, out_ch, kernel_size=3, padding=1),
        nn.BatchNorm2d(out_ch),
        nn.ReLU(inplace=True),
    )
\end{lstlisting}

Cách thứ hai là nạp mã từ một file có sẵn trong đồ án. Cách này tiện vì báo
cáo luôn khớp với mã thật, sửa mã là báo cáo tự cập nhật theo. Mã
nguồn~\ref{lst:tufile} được nạp từ file \code{code/vi-du-huan-luyen.py}.

\lstinputlisting[language=Python, caption={Chèn mã nguồn từ file có sẵn.}, label={lst:tufile}]{code/vi-du-huan-luyen.py}

\begin{luuy}
Quy ước của khoa: comment trong code viết bằng tiếng Anh, còn văn xuôi của
báo cáo viết bằng tiếng Việt. Không đặt tiếng Việt bên trong khối mã nguồn,
vì phông mã nguồn không có đủ dấu tiếng Việt.
\end{luuy}

Chỉ đưa vào thân bài những đoạn mã cốt lõi. Mã dài để ở phụ lục hoặc chỉ dẫn
tới kho \code{git} của đồ án.

\section{Biểu diễn giải thuật}

Giải thuật trình bày bằng mã giả, không phải mã của một ngôn ngữ cụ thể,
để người đọc nắm được ý tưởng mà không cần biết ngôn ngữ đó.
Giải thuật~\ref{alg:nguyento} là một ví dụ.

\begin{algorithm}[H]
\caption{Kiểm tra một số tự nhiên có phải số nguyên tố hay không}
\label{alg:nguyento}
\hspace*{\algorithmicindent}\textbf{Đầu vào:} số tự nhiên $n$\\
\hspace*{\algorithmicindent}\textbf{Đầu ra:} \textsc{True} nếu $n$ là số nguyên tố, ngược lại \textsc{False}
\begin{algorithmic}[1]
\Function{LaSoNguyenTo}{$n$}
  \If{$n < 2$}
    \State \Return \textsc{False}
  \EndIf
  \For{$i \gets 2$ \textbf{to} $\lfloor \sqrt{n} \rfloor$}
    \If{$n \bmod i = 0$}
      \State \Return \textsc{False}
    \EndIf
  \EndFor
  \State \Return \textsc{True}
\EndFunction
\end{algorithmic}
\end{algorithm}

\section{Trích dẫn tài liệu tham khảo}

Tài liệu tham khảo khai báo trong file \code{refs.bib}, mỗi tài liệu một
khối. Trong thân bài em gọi bằng lệnh \code{cite} kèm khoá của tài liệu, ví
dụ mô hình AlexNet~\cite{Krizhevsky2012} và mạng ResNet~\cite{He2016}. LaTeX
tự đánh số và tự dựng danh mục ở cuối báo cáo, nên em không phải đánh số tay.

\begin{ghinho}
Chỉ liệt kê trong \code{refs.bib} những tài liệu em thực sự trích dẫn. Tài
liệu khai báo mà không gọi \code{cite} sẽ không xuất hiện trong danh mục.
\end{ghinho}

\section{Chú thích cuối trang}

Khi cần giải thích thêm mà không muốn làm gãy mạch câu, dùng lệnh
\code{footnote}\footnote{Đây là một dòng chú thích cuối trang.}. Đừng lạm
dụng, mỗi trang nhiều lắm một đến hai chú thích.

\section{Các hộp nhấn mạnh}

Mẫu này có sẵn bốn loại hộp để làm nổi những ý quan trọng.

\begin{dinhnghia}[Học chuyển giao]
Học chuyển giao \en{(transfer learning)} là kỹ thuật tận dụng tri thức mà mô
hình đã học được trên một bài toán nguồn có nhiều dữ liệu, để giải một bài
toán đích có ít dữ liệu hơn.
\end{dinhnghia}

\begin{ghinho}
Hộp ghi nhớ dùng cho ý chốt mà em muốn người đọc mang theo sau khi đọc xong
mục đó.
\end{ghinho}

\begin{luuy}
Hộp lưu ý dùng cho cảnh báo, điều kiện áp dụng, hoặc lỗi mà người đọc dễ mắc
nếu làm theo.
\end{luuy}

\begin{vidu}[minh hoạ một tình huống cụ thể]
Hộp ví dụ tự đánh số theo chương, tiện khi em cần dẫn lại một ví dụ ở chỗ
khác trong báo cáo.
\end{vidu}


% ---------------------------------------------------------------------
%  TÀI LIỆU THAM KHẢO VÀ PHỤ LỤC
% ---------------------------------------------------------------------
\backmatter

% \phantomsection tạo mốc neo riêng cho mục lục, nếu thiếu thì hyperref
% báo trùng neo giữa dấu trang và dấu trang cha.
\cleardoublepage
\phantomsection
\addcontentsline{toc}{chapter}{Tài liệu tham khảo}
\bibliographystyle{plain}
\bibliography{refs}

%% ---------------------------------------------------------------------
%  PHỤ LỤC
% ---------------------------------------------------------------------
% Trong \backmatter thì \chapter đã tự thêm vào mục lục, không cần
% \addcontentsline nữa, thêm vào là trùng neo dấu trang.
\chapter{Phụ lục}
\markboth{Phụ lục}{}

\begin{chuy}
\goiy{phụ lục là nơi đặt những nội dung dài không tiện chèn vào thân báo
cáo: mã nguồn đầy đủ, bảng dữ liệu mở rộng, kết quả thử nghiệm bổ sung, biên
bản hướng dẫn. Riêng mục A và mục B nên giữ lại.}
\end{chuy}

\section*{A. Hướng dẫn cài đặt và chạy chương trình}
\phantomsection
\addcontentsline{toc}{section}{A. Hướng dẫn cài đặt và chạy chương trình}

Viết các bước để một người chưa từng thấy đồ án vẫn chạy được sản phẩm. Hội
đồng sẽ chạy thử theo đúng các bước này.

\begin{lstlisting}[language=bash, caption={Các bước cài đặt và khởi chạy ứng dụng.}]
# 1. Clone the repository
git clone https://github.com/tencuaem/ten-do-an.git
cd ten-do-an

# 2. Create a virtual environment and install dependencies
python -m venv .venv
.venv\Scripts\activate
pip install -r requirements.txt

# 3. Train the model (skip if a trained checkpoint is provided)
python train.py --epochs 30 --data data/

# 4. Launch the web application, then open http://localhost:5000
python app.py
\end{lstlisting}

\section*{B. Bảng tự kiểm trước khi nộp}
\phantomsection
\addcontentsline{toc}{section}{B. Bảng tự kiểm trước khi nộp}

Trước khi in quyển, em tự đánh dấu vào từng ô dưới đây.

\begin{tukiem}
  \item Sản phẩm chạy được trên máy khác, không chỉ trên máy của em.
  \item Đã chạy thử demo trên đúng môi trường sẽ dùng khi bảo vệ.
  \item Mã nguồn nằm trên kho \code{git}, lịch sử commit cho thấy tiến độ
        trải đều chứ không dồn vào tuần cuối.
  \item Có tệp \code{README} và \code{requirements.txt} để người khác chạy lại.
  \item Mọi hình, bảng và giải thuật đều có chú thích, có nhãn, và được nhắc
        tới trong thân bài.
  \item Mọi mục tiêu ở Chương~\ref{ch:modau} đều được đối chiếu ở
        Chương~\ref{ch:ketluan}.
  \item Mọi con số trong báo cáo trùng khớp với kết quả chạy thật, không có
        số nào là ước lượng hay bịa.
  \item Mọi tài liệu trong \code{refs.bib} đều được trích dẫn ít nhất một
        lần, và mọi nguồn dữ liệu, mã tham khảo đều đã ghi rõ.
  \item Đã khai báo phần nào có dùng công cụ trí tuệ nhân tạo hỗ trợ, và em
        giải thích được toàn bộ hệ thống khi bị hỏi.
  \item Đã xoá hết hộp gợi ý, mọi khối nội dung mẫu, và cả
        Chương~\ref{ch:soanthao} vốn chỉ là sổ tra cứu.
\end{tukiem}

\begin{luuy}
Hai mục cuối là lỗi hay gặp nhất khi dùng mẫu có sẵn. Trước khi nộp hãy tìm
trong bản PDF xem còn dòng chữ xám \emph{NỘI DUNG MẪU} nào sót lại không.
\end{luuy}

\section*{C. Nhận xét của giảng viên hướng dẫn}
\phantomsection
\addcontentsline{toc}{section}{C. Nhận xét của giảng viên hướng dẫn}

\khungnhanxet[7.5cm]

\vspace{0.8em}
\noindent\begin{minipage}[t]{0.5\textwidth}
\color{daunavy}
{\color{daumaroon}\bfseries Điểm số}

\vspace{0.7em}
\dongchamnhan{Bằng số:}

\vspace{0.9em}
\dongchamnhan{Bằng chữ:}
\end{minipage}\hfill
\okyten{\dauGVHD}{Giảng viên hướng dẫn}


\end{document}
